\documentclass[12pt]{article}

\usepackage{graphicx}
\usepackage{amsmath}
\usepackage{amsfonts}
\usepackage{amssymb}
\usepackage{enumitem}
\usepackage{hyperref}
\usepackage{geometry}
\geometry{margin=1in}

\title{BMS Firmware Development Plan}
\author{Kaushik Vada}
\date{}

\begin{document}

\maketitle
\tableofcontents
\newpage

\section{Project Overview}

The battery management system utilizes a 10s1p Molicel P42A pack ($\sim$36~V, 4.2~Ah) instrumented with NTC thermistors and supported by an active fan cooling system. The isolated electronics block contains an STM32 microcontroller, an 8-channel isolated SPI ADC, an I2C real-time clock, isolated DC/DC supplies for 5V and 3.3V, and a USB-C CDC link to a host computer. Safety mechanisms include a PTC fuse and comparator-driven MOSFET gate disabling.

The firmware must periodically acquire voltage, current, and temperature data, control cooling and protection mechanisms, ensure fail-safe operation, and stream structured telemetry to a host graphical interface.

\section{Block Diagram Summary}

\subsection{Battery Front-End}

\begin{itemize}
\item 10s1p Molicel P42A cell group with low-side shunt current sensing.
\item BQ76930 BMS IC providing voltage sensing, current reporting, cell balancing, and fault reporting.
\item Low-side electronic breaker (with redundancy in revision 2).
\item Per-cell balance resistors and capacitance.
\item Optional comparator and discrete gate driver for redundant hardware shutdown.
\end{itemize}

\subsection{Thermal Sensing}

\begin{itemize}
\item Ten NTC thermistors routed to an isolated ADC.
\item Divider network and low-pass filters referenced to \texttt{VREF\_THERM}.
\item Temperature sampling rate of 5--10 Hz.
\item Exceeding thermal threshold triggers GUI alert, LED indicator, and breaker opening.
\end{itemize}

\subsection{Power and Isolation}

\begin{itemize}
\item Always-on power path ensures STM32 remains active even when breaker opens.
\item PTC fuse protects isolated 5V rail.
\item 5V-to-3.3V isolated DC/DC regulation.
\item ESD protection, watchdog, and supervisor ICs included in revision 2.
\end{itemize}

\subsection{Control Outputs}

\begin{itemize}
\item PWM fan output routed through isolated driver/optocoupler.
\item Comparator outputs driving MOSFET gate disable circuitry.
\item Redundant breaker-actuation GPIO.
\item Isolated I2C connection to real-time clock.
\end{itemize}

\subsection{Host Interface and Testbench Setup}

\begin{itemize}
\item USB-C CDC link to desktop GUI.
\item Bench testing includes charger and electronic load.
\item Optional NVIDIA Jetson downstream consumer for ML telemetry.
\end{itemize}

\section{Firmware Architecture}

\subsection{Required Drivers}

\begin{itemize}
\item I2C for BQ76930 battery front end.
\item SPI or I2C for isolated 10-channel ADC.
\item I2C connection to real-time clock.
\item PWM timers for cooling fan(s).
\item GPIO for breaker, fault inputs, and comparator signals.
\item USB-CDC for structured host communication.
\item Windowed watchdog for safety and firmware reliability.
\end{itemize}

\subsection{Core State Machines}

\subsubsection{Temperature State Machine}

\begin{itemize}
\item Sample NTC array at 5--10 Hz.
\item Compare temperatures against warning and cutoff thresholds.
\item On trip: log event, trigger USB alert, activate LED, and open breaker.
\end{itemize}

\subsubsection{State of Charge (SOC) Machine}

\begin{itemize}
\item Integrate shunt current using BQ76930 coulomb-count reporting.
\item Reset logic triggered on charge-source plug/unplug.
\item SOC bounds limit system behavior and GUI reporting.
\end{itemize}

\subsubsection{Safety Path}

\begin{itemize}
\item BQ76930 interrupts and comparator signals serviced by highest priority ISRs.
\item Hardware trip opens breaker regardless of firmware execution state.
\item MCU remains powered to retrieve and transmit failure information.
\end{itemize}

\subsection{Communications Protocol}

\begin{itemize}
\item USB-CDC transmits structured frames containing:
\begin{itemize}
\item Timestamp
\item Pack voltage and current
\item Per-cell voltages and temperatures
\item Fan RPM
\item Operating mode
\item Fault bitmap (OV, UV, OT, OCP, comparator trip)
\end{itemize}
\item Accept inbound control messages (fan override, mode selection, balancing enable).
\end{itemize}

\section{GUI Integration Strategy}

\subsection{Expected JSON Schema}

\begin{verbatim}
{
cells:[{id, voltage, temperature}],
mode:"Balanced",
fan1:{rpm},
fan2:{rpm}
}
\end{verbatim}

\subsection{Python Host Bridge}

\begin{itemize}
\item Implement \texttt{backend/data\_stream.py}.
\item Parse binary USB-CDC telemetry frames.
\item Broadcast live JSON via WebSocket at:
[
\texttt{ws://localhost:8766/telemetry}
]
\item Relay GUI-originating commands to MCU.
\end{itemize}

\subsection{Frontend Modifications}

\begin{itemize}
\item Replace mock timers in \texttt{scene.js} with WebSocket consumer.
\item Update GUI rendering functions on message receipt.
\item Maintain fallback mock data if socket disconnects.
\end{itemize}

\subsection{Application Launcher Alignment}

\begin{itemize}
\item \texttt{gui\_launcher.py} already serves application assets.
\item Automatically launch Python bridge for turnkey execution.
\end{itemize}

\section{Firmware Bring-Up Procedure}

\subsection{STM32CubeMX Project Setup}

\begin{itemize}
\item Finalize microcontroller selection and pinout.
\item Enable USB-CDC, I2C peripherals, SPI/I2C ADC interface, PWM timers, and EXTI lines.
\item Configure 1 kHz system tick; RTOS optional.
\end{itemize}

\subsection{Sensor Acquisition Layer}

\begin{itemize}
\item Perform DMA-based SPI acquisition at 10--20 Hz.
\item Apply calibration, scaling, and noise filtering.
\item Timestamp packets using RTC.
\end{itemize}

\subsection{Control and Safety Logic}

\begin{itemize}
\item Implement voltage, temperature, and current protection thresholds.
\item Drive MOSFET enable logic according to comparator state.
\item Execute fan control using threshold or PID strategies.
\item Verify isolated power-good before enabling communication.
\end{itemize}

\subsection{USB Frame Format}

\begin{itemize}
\item Header, payload length, data structure, CRC.
\item Maintain ring buffer to prevent USB-induced blocking.
\end{itemize}

\subsection{Initial Bench Validation}

\begin{itemize}
\item Validate ADC readout on scope and against multimeter.
\item Test RTC communication over I2C.
\item Sweep fan PWM duty cycle and verify RPM feedback.
\item Confirm fast hardware cutoff via comparator.
\item Stream test text/JSON over USB for GUI debugging.
\end{itemize}

\section{End-to-End System Validation}

\begin{enumerate}
\item Finalize USB and WebSocket message formats.
\item Implement Python bridge with synthetic/mock frames.
\item Verify GUI displays live telemetry.
\item Replace mock data with real firmware output.
\item Perform combined load, charge, and thermal testing.
\item Integrate with full pack after bench-level signoff.
\end{enumerate}

\section{Recommended Next Steps}

\begin{enumerate}
\item Lock down MCU selection and hardware pin mapping.
\item Generate and version-control CubeMX firmware skeleton.
\item Implement firmware telemetry framing and decoding.
\item Build Python WebSocket bridge with replayable mock data.
\item Connect GUI and verify real-time responsiveness.
\item Conduct hardware-in-the-loop testing.
\end{enumerate}

\end{document}
